\documentclass[12pt]{article}
\usepackage[utf8]{inputenc}
\usepackage[margin=1in]{geometry}
\usepackage{amsmath}
\usepackage{amssymb}
\usepackage{natbib}
\usepackage{hyperref}
\usepackage{booktabs}
\usepackage{graphicx}
\usepackage{setspace}
\onehalfspacing

\title{Land Speculation, Public Debt, and Hamilton's Funding System:\\
The Political Economy of Competing Interests, 1783--1795}

\author{Prepared for Thomas J. Sargent and George J. Hall}

\date{\today}

\begin{document}

\maketitle

\begin{abstract}
This paper examines the conflict between land speculators and public creditors in the 1780s and analyzes how Alexander Hamilton's assumption and funding program (1790--1791) resolved this conflict by eliminating the arbitrage opportunities that had made depreciated Continental securities valuable for purchasing western lands. Drawing on Ferguson's analysis of state-building through fiscal policy and Rothbard's emphasis on special interest conflicts, we document the mechanisms of land-for-debt schemes, identify key speculators, quantify the magnitudes involved, and trace the consequences of Hamilton's measures for different interest groups. The episode illustrates how fiscal policy choices can fundamentally restructure property rights and political coalitions.
\end{abstract}

\section{Introduction}

The fiscal history of the United States between the Treaty of Paris (1783) and Hamilton's funding system (1790--1791) reveals a fundamental conflict of interests that has received insufficient attention in standard treatments of the early Republic. While historians have long understood that Hamilton's assumption and funding of Revolutionary War debts represented a crucial act of state-building \citep{Ferguson1961}, the specific mechanism by which these measures affected competing economic interests deserves closer analysis.

Several states in the 1780s established schemes whereby western public lands could be purchased using Continental securities or state debt certificates at face value, even though these securities traded at substantial discounts---typically 15--25 cents on the dollar---in secondary markets. This created enormous arbitrage opportunities for speculators with capital and information \citep{Rothbard1999, Rothbard2011}. Prominent figures including Robert Morris, William Duer, Oliver Phelps, Nathaniel Gorham, and many others accumulated vast land holdings using depreciated paper.

The political economy of this situation was complex. Land speculators needed depreciated securities to make their schemes profitable. Public creditors---both domestic and foreign---wanted securities to appreciate toward par. States with western lands could retire debts cheaply through land sales. States without western lands sought federal assumption of debts for equity reasons. Hamilton's program resolved this conflict decisively in favor of public creditors and federal power, but at the cost of bankrupting many prominent speculators and ending states' fiscal autonomy.

This paper proceeds as follows. Section~\ref{sec:debt} describes the structure and market value of Continental and state debts in the 1780s. Section~\ref{sec:schemes} documents the major land speculation schemes and identifies key players. Section~\ref{sec:conflict} analyzes the conflict of interests between speculators and creditors. Section~\ref{sec:hamilton} explains Hamilton's measures and their passage. Section~\ref{sec:effects} quantifies the distributional effects on different groups. Section~\ref{sec:consequences} discusses long-run consequences for American fiscal development. Section~\ref{sec:conclusion} concludes.

\section{The Structure of Revolutionary War Debts}\label{sec:debt}

\subsection{Continental Debt}

The Continental Congress emerged from the Revolutionary War with approximately \$42 million in domestic debt, composed of several types of instruments:

\begin{itemize}
\item \textbf{Loan Office Certificates} ($\sim$\$11 million): Issued to lenders during the war, bearing 6\% interest
\item \textbf{Army Debt Certificates} ($\sim$\$10 million): Pay certificates issued to soldiers and officers
\item \textbf{Quartermaster and Commissary Certificates} ($\sim$\$16 million): Issued for supplies requisitioned during the war
\item \textbf{Final Settlement Certificates} ($\sim$\$5 million): Issued to settle accounts after the war
\end{itemize}

In addition, the Continental Congress owed approximately \$12 million to foreign creditors, primarily France and Dutch banking houses. This foreign debt was in a separate category because the Confederation Congress, despite its weakness, maintained service on foreign obligations to preserve access to European credit markets.

\subsection{State Debts}

The thirteen states collectively held approximately \$25 million in war-related debts. These varied substantially across states:

\begin{itemize}
\item \textbf{Massachusetts}: $\sim$\$5 million (had been servicing debt through heavy taxation)
\item \textbf{South Carolina}: $\sim$\$5 million (heavy war damage, planter debt)
\item \textbf{Virginia}: $\sim$\$3.5 million (but had already paid much more)
\item \textbf{Connecticut}: $\sim$\$1.6 million
\item Other states: smaller amounts
\end{itemize}

The key distinction was between states that had heavily taxed themselves to service debts (notably Massachusetts, leading to Shays' Rebellion in 1786--87) and states that had done little (South Carolina, North Carolina, Georgia).

\subsection{Secondary Market Prices}

Continental securities traded at severe discounts in the 1780s:

\begin{center}
\begin{tabular}{ll}
\toprule
Period & Market Price (cents on dollar) \\
\midrule
1784--1786 & 15--20 \\
1787--1788 & 20--25 \\
1789 & Rising gradually (25--40) \\
1790 (pre-funding) & 40--50 \\
1791 (post-funding) & 80--100 \\
\bottomrule
\end{tabular}
\end{center}

These discounts reflected several factors:
\begin{enumerate}
\item Uncertainty about the federal government's ability and willingness to pay
\item Lack of federal taxing power under the Articles of Confederation
\item States' varying commitment to servicing their own debts  
\item General economic depression following the war
\item Expected inflation given Congress's history of currency emissions
\end{enumerate}

State securities traded at comparable or even steeper discounts, with substantial variation across states reflecting their differing fiscal conditions.

\section{The Land Speculation Schemes}\label{sec:schemes}

The depreciation of public securities created arbitrage opportunities when states made western lands available for purchase with these securities accepted at face value. We document the major schemes:

\subsection{North Carolina Western Lands}

North Carolina's 1783 land law allowed purchase of western lands (in what would become Tennessee) with depreciated Continental certificates and state paper at face value. Key speculators included:

\begin{itemize}
\item \textbf{William Blount} (1749--1800): Accumulated hundreds of thousands of acres; later served as governor of the Southwest Territory and as Senator from Tennessee
\item \textbf{John Sevier} (1745--1815): Revolutionary War hero who became a massive land speculator and first governor of Tennessee
\end{itemize}

The arithmetic was compelling: a speculator purchasing \$1,000 face value of Continental certificates for \$200 (20 cents on the dollar) could then acquire 200,000 acres at the official rate of approximately 10 shillings per 100 acres.

\subsection{The Georgia Yazoo Schemes}

Georgia's western land schemes (covering present-day Alabama and Mississippi) became notorious:

\subsubsection{The First Yazoo Sale (1789)}

Three companies (Virginia Company, South Carolina Company, Tennessee Company) purchased approximately 25.4 million acres for roughly \$200,000, payable in depreciated state certificates. Many Georgia legislators received bribes. The deal collapsed amid scandal, but it illustrated the scale of speculation.

\subsubsection{The Second Yazoo Sale (1795)}

Four companies purchased approximately 35 million acres for \$500,000. Virtually every Georgia legislator received bribes or shares. The sale was eventually overturned by a subsequent legislature, leading to decades of litigation and eventually the landmark Supreme Court case \textit{Fletcher v. Peck} (1810), which established that legislative rescission of a contract violated the Constitution's Contract Clause.

Prominent Yazoo speculators included:
\begin{itemize}
\item Patrick Henry: Virginia's great Revolutionary orator
\item Wade Hampton I: South Carolina planter and Revolutionary War officer
\item Robert Morris: through various intermediaries
\item Numerous sitting members of Congress
\end{itemize}

\subsection{The Phelps-Gorham Purchase (1788)}

This Massachusetts transaction was perhaps the best-documented major land speculation of the era:

\textbf{The Transaction:}
\begin{itemize}
\item Oliver Phelps (1749--1809) and Nathaniel Gorham (1738--1796) purchased approximately 2.6 million acres in western New York from Massachusetts
\item Nominal price: \$1 million in Massachusetts securities
\item Actual payment: approximately \$300,000 in depreciated Massachusetts paper (at $\sim$21 cents on the dollar)
\item Expected resale at much higher prices to settlers
\end{itemize}

\textbf{The Collapse:}
\begin{itemize}
\item Hamilton's funding plan passed (1790)
\item Massachusetts securities appreciated rapidly toward par
\item Phelps and Gorham still owed installment payments totaling \$700,000
\item In appreciated securities, they now effectively owed \$1 million rather than \$200,000
\item Unable to meet payments, they forfeited most of the land back to Massachusetts in 1790
\item Gorham, who had been a delegate to the Constitutional Convention, died in debtor's circumstances in 1796
\end{itemize}

This case perfectly illustrates how Hamilton's funding destroyed the land speculation arbitrage.

\subsection{The Robert Morris Empire}

Robert Morris (1734--1806)---``Financier of the Revolution,'' delegate to the Continental Congress, Senator from Pennsylvania---built the most extensive land speculation empire:

\textbf{Holdings (1790s):}
\begin{itemize}
\item Lots in the new federal capital (Washington, DC)
\item Approximately 6 million acres across Pennsylvania, Virginia, North Carolina, Georgia, and South Carolina
\item Partner with John Nicholson (Pennsylvania's former Comptroller-General) and James Greenleaf in numerous ventures
\end{itemize}

\textbf{The Financial Logic:}
Morris's calculations assumed:
\begin{enumerate}
\item Continued availability of cheap Continental paper for land purchases
\item Steady appreciation of land values as settlement progressed
\item Ability to resell parcels at profits sufficient to service acquisition debt
\end{enumerate}

Hamilton's funding destroyed assumption (1). The depression of the mid-1790s destroyed assumption (2). By 1798, Morris was imprisoned in Philadelphia's Prune Street Prison for debt, remaining there until passage of the federal bankruptcy law in 1801.

\subsection{Other Major Schemes}

\textbf{The Symmes Purchase (Ohio):} Judge John Cleves Symmes (1742--1814) purchased approximately 1 million acres between the Miami Rivers in Ohio (1788) using depreciated Continental certificates. He failed to complete payments when securities appreciated and lost most of his claim. The town of Cincinnati developed within his original purchase area.

\textbf{New York State Lands:} Beyond the Phelps-Gorham purchase, New York made vast tracts available for purchase with Continental securities, attracting numerous speculators.

\textbf{Southern Schemes:} South Carolina and Virginia also had land-for-debt provisions, though smaller in scale than North Carolina and Georgia.

\section{The Conflict of Interests}\label{sec:conflict}

\subsection{Land Speculators' Position}

For land speculators, the arbitrage opportunity was clear and immediate:
\begin{enumerate}
\item Purchase Continental or state securities at market prices (15--25 cents on dollar)
\item Use these securities at face value to purchase western lands from states
\item Resell land parcels to settlers at market prices
\item Expected return: 300--500\% or more
\end{enumerate}

\textbf{Crucially, this strategy required depreciated securities.} If Continental securities rose to par, the arbitrage disappeared entirely. A speculator who had already committed to land purchases at par value (using securities purchased at 20 cents) would face disaster if securities appreciated before the transaction closed.

This created a perverse incentive structure: sophisticated speculators---many of whom held large portfolios of Continental securities---had an interest in \textit{preventing} those securities from appreciating.

\subsection{Public Creditors' Position}

Public creditors wanted exactly the opposite. This group included:

\begin{enumerate}
\item \textbf{Foreign creditors:} French government and Dutch banking houses (approximately \$12 million)
\item \textbf{Domestic creditors:} By 1790, Continental debt had become concentrated in commercial centers through market transactions:
\begin{itemize}
\item Pennsylvania: $\sim$\$8 million (especially Philadelphia merchants)
\item Massachusetts: $\sim$\$7 million (especially Boston merchants)  
\item New York: $\sim$\$5 million (especially New York City merchants)
\item Maryland: $\sim$\$3 million
\item South Carolina: $\sim$\$3 million
\item Other states: smaller amounts
\end{itemize}
\item \textbf{Original holders:} Most soldiers, farmers, and widows who received certificates had already sold them at steep discounts, so by 1790 relatively few original holders remained
\end{enumerate}

These creditors wanted:
\begin{itemize}
\item A credible federal funding plan
\item Regular interest payments
\item Eventual principal repayment or at least maintenance of principal value
\item Liquid, negotiable securities that could serve as money-like assets
\end{itemize}

\subsection{States' Divergent Interests}

States faced conflicting incentives based on their circumstances:

\textbf{States with Western Land Claims} (Virginia, North Carolina, South Carolina, Georgia, New York):
\begin{itemize}
\item Could use land-for-debt schemes to retire obligations cheaply
\item Resisted ceding land claims to federal government
\item Saw land sales as major revenue source
\item Generally opposed federal assumption (preferred state control)
\end{itemize}

\textbf{States without Western Land Claims} (Maryland, New Jersey, Delaware, Rhode Island):
\begin{itemize}
\item Lacked land-for-debt option
\item Needed other revenue sources for debt service
\item Supported federal assumption for equity reasons
\item Wanted compensation for states that had contributed land and resources
\end{itemize}

\textbf{States that had Already Paid Their Debts} (Massachusetts, partially Virginia):
\begin{itemize}
\item Had heavily taxed citizens to service war debts
\item Wanted federal assumption to provide tax relief
\item Felt equity demanded sharing the burden across all states
\item Shays' Rebellion (1786--87) in Massachusetts had been partly triggered by these taxes
\end{itemize}

\subsection{The Option Value Calculation}

William Duer's position illustrates the complexity. As Assistant Secretary of the Treasury (1789--90), Duer had inside knowledge of Hamilton's funding plans. Yet he pursued both strategies simultaneously:

\begin{enumerate}
\item Purchased Continental securities on margin, speculating on appreciation
\item Used securities and borrowed funds to purchase western lands
\item Expected to profit from both strategies
\end{enumerate}

The leverage destroyed him. When the securities market broke in March 1792, he could not meet margin calls. His bankruptcy triggered the Panic of 1792. He spent the rest of his life (until death in 1799) in New York's debtor's prison.

Duer's failure illustrates that the two uses of Continental securities---as appreciating financial assets versus discounted land currency---were fundamentally incompatible. A speculator had to choose one strategy or the other; attempting both with leverage was catastrophic.

\section{Hamilton's Measures}\label{sec:hamilton}

\subsection{The First Report on Public Credit (January 1790)}

Alexander Hamilton's \textit{First Report on the Public Credit} laid out a comprehensive plan \citep{Hamilton1790}:

\textbf{Key Provisions:}
\begin{enumerate}
\item \textbf{Fund all Continental debt at par:} Approximately \$42 million to be exchanged for new federal bonds
\item \textbf{Assume state debts:} Federal government to assume approximately \$25 million in state war debts
\item \textbf{Create sinking fund:} Dedicate specific revenues to gradual debt reduction
\item \textbf{Establish permanent revenues:} Pledge tariff revenues and new excise taxes to debt service
\end{enumerate}

\textbf{The Funding Terms (Menu of Options):}

Hamilton offered holders of old Continental securities several choices:

\begin{itemize}
\item \textbf{Option 1:} New bonds paying 6\% annually in perpetuity, but with face value reduced to approximately 2/3 of original

\item \textbf{Option 2:} New bonds at full face value paying:
\begin{itemize}
\item 6\% on 2/3 of principal (immediate)
\item 6\% on 1/3 of principal (deferred 10 years)
\end{itemize}

\item \textbf{Option 3:} New bonds at face value paying 4\% in perpetuity, plus claims to western lands
\end{itemize}

Most holders chose Option 2, accepting an effective yield of approximately 4--4.5\% (given the present value of deferred interest) in exchange for certainty of payment.

\textbf{Revenue Sources:}
\begin{itemize}
\item Existing tariffs (now permanently pledged to debt service)
\item New excise taxes, especially on distilled spirits (leading to the Whiskey Rebellion of 1794)
\item Future land sales from federal (not state) western lands
\end{itemize}

\subsection{The Political Battle}

\subsubsection{Madison's Objections}

James Madison, Hamilton's former ally in the Federalist Papers, raised several objections in Congress:

\begin{enumerate}
\item \textbf{The ``Original Holders'' Argument:} Most Continental debt had transferred from soldiers, farmers, and widows who sold at steep discounts to speculators. Funding at par would reward speculators rather than patriots. Madison proposed ``discrimination''---paying current holders market value and compensating original holders for the difference. This was rejected as:
\begin{itemize}
\item Administratively impractical (certificates had changed hands many times)
\item Destructive to negotiability of government securities
\item Based on incomplete records
\end{itemize}

\item \textbf{Regional Equity:} Southern states (particularly Virginia) had already paid off much of their war debt through heavy taxation. Assumption would require them to help pay northern states' debts through federal taxes. This seemed inequitable.

\item \textbf{Constitutional Concerns:} Madison questioned whether the federal government possessed constitutional authority to assume state debts, given that these obligations had not been contracted by the federal government.
\end{enumerate}

\subsubsection{Hamilton's Defense}

Hamilton responded on multiple grounds:

\begin{enumerate}
\item \textbf{Practical Necessity:} Discrimination between original and current holders was impossible given poor records and multiple transfers. More fundamentally, it would destroy the negotiability of government securities---no one would accept government paper if payment could later be challenged based on transfer history.

\item \textbf{Economic Benefits:} A funded national debt would:
\begin{itemize}
\item Establish public credit, lowering interest rates for all borrowers
\item Create liquid securities that could circulate as money-like assets
\item Facilitate commerce and economic development
\item Enable future government borrowing for national emergencies
\item Stimulate securities markets and financial sector development
\end{itemize}

\item \textbf{Political Necessity:} Hamilton argued that a funded national debt would create a powerful interest group committed to the federal government's success and stability. This was exactly what the weak Confederation had lacked. Binding creditor interests to national government would strengthen the union.

\item \textbf{Justice:} All states had benefited from the war. All should share the costs. States that had paid more should not be penalized; rather, assumption would provide them tax relief going forward.
\end{enumerate}

\subsubsection{The Compromise of 1790}

The famous ``Dinner Table Bargain'' (June 1790) involved Jefferson, Hamilton, and Madison:
\begin{itemize}
\item Hamilton obtained votes for assumption of state debts
\item Virginia delegation obtained the permanent federal capital on the Potomac River
\item Pennsylvania obtained the temporary capital in Philadelphia (1790--1800)
\end{itemize}

However, this understates the complexity of vote-trading. Hamilton made specific side deals:
\begin{itemize}
\item \textbf{Maryland:} Federal assumption of Maryland's debt plus capital location near Maryland
\item \textbf{Pennsylvania:} Assumption plus temporary capital
\item \textbf{Individual Congressmen:} Various adjustments to their states' assumed debt amounts
\end{itemize}

\textbf{Final Votes (July 1790):}
\begin{itemize}
\item House of Representatives: 34--28 in favor of assumption
\item Senate: Passed more easily
\end{itemize}

\subsection{The Report on a National Bank (December 1790)}

Hamilton's \textit{Report on a National Bank} complemented the funding system \citep{Hamilton1790bank}:

\textbf{Provisions:}
\begin{itemize}
\item Create Bank of the United States with \$10 million capital
\item Private subscription of \$8 million; federal government subscription of \$2 million
\item Accept government securities for three-quarters of private subscription
\item 20-year charter
\item Bank authorized to establish branches
\item Bank notes to be accepted for tax payments
\end{itemize}

\textbf{Functions:}
\begin{itemize}
\item Facilitate government borrowing and debt management
\item Provide short-term credit to Treasury
\item Create a uniform national currency (bank notes)
\item Serve as a lender to the economy
\item Hold government deposits
\end{itemize}

The Bank was established in 1791 and its stock subscription was heavily oversubscribed, demonstrating the success of the funding system in creating a class of investors with capital and confidence in federal securities.

\subsection{Total Funded Debt (1791)}

After Hamilton's measures:
\begin{itemize}
\item Continental debt funded: $\sim$\$42 million
\item State debts assumed: $\sim$\$25 million
\item Foreign debt (separately arranged): $\sim$\$12 million
\item \textbf{Total funded federal debt: $\sim$\$79 million}
\end{itemize}

This represented approximately 40\% of estimated GDP at the time---a substantial but manageable debt burden given the dedicated revenue sources.

\section{Distributional Effects}\label{sec:effects}

\subsection{Effects on Bondholders}

Public creditors experienced substantial windfall gains. Consider a representative Philadelphia merchant:

\begin{center}
\begin{tabular}{lll}
\toprule
Year & Face Value & Market Value \\
\midrule
1786 (purchase) & \$10,000 & \$2,000 \\
1789 (anticipation) & \$10,000 & \$4,000 \\
1790 (mid-year) & \$10,000 & \$5,000--6,000 \\
1791 (after funding) & \$10,000 & \$8,000--9,000 \\
\bottomrule
\end{tabular}
\end{center}

This merchant realized an annual return of approximately 35--40\% over five years. Moreover, the new federal bonds paid regular interest (4--6\% depending on option chosen), providing ongoing income.

\textbf{Geographic Distribution of Gains:}

Ferguson's research documented that Continental debt had become heavily concentrated in northern commercial centers by 1790 \citep{Ferguson1961}:

\begin{itemize}
\item Pennsylvania (especially Philadelphia): Held $\sim$\$8 million, gained $\sim$\$5--6 million in wealth from appreciation
\item Massachusetts (especially Boston): Held $\sim$\$7 million, gained $\sim$\$4--5 million
\item New York (especially New York City): Held $\sim$\$5 million, gained $\sim$\$3--4 million
\end{itemize}

These gains accrued primarily to merchants, not to original certificate holders (soldiers, farmers, widows) who had sold years earlier.

\subsection{Effects on Land Speculators}

\subsubsection{The Phelps-Gorham Disaster}

The Phelps-Gorham case illustrates the arithmetic:

\begin{itemize}
\item \textbf{1788:} Purchased 2.6 million acres for \$1 million in Massachusetts securities
\item \textbf{1788:} Paid approximately \$300,000 in depreciated paper (30\% down payment at $\sim$20 cents on dollar)
\item \textbf{Outstanding obligation:} \$700,000 in future installments
\item \textbf{1790:} Massachusetts securities rose from 20 cents to near par
\item \textbf{Effective cost:} Jumped from expected \$300,000 total to actual \$1,000,000
\item \textbf{Problem:} Had already resold some parcels based on the old, lower cost basis
\item \textbf{Result:} Unable to meet installment payments; forfeited most of the purchase back to Massachusetts in 1790
\end{itemize}

Gorham died in straitened circumstances in 1796. Phelps survived but never recovered financially.

\subsubsection{The Morris-Nicholson-Greenleaf Collapse}

The Morris-Nicholson-Greenleaf partnership assembled approximately 6 million acres in the 1790s through various schemes. Much was purchased with borrowed money, with plans to repay loans through land sales. The depression of 1796--97 crushed land values. Results:

\begin{itemize}
\item \textbf{Robert Morris:} Imprisoned for debt 1798--1801 in Philadelphia's Prune Street Prison; released only after passage of federal bankruptcy law
\item \textbf{John Nicholson:} Died in Philadelphia's debtor's prison in 1800
\item \textbf{James Greenleaf:} Bankrupted; eventually recovered partially
\end{itemize}

Their creditors eventually recovered some amounts through long legal processes, but the principals lost everything.

\subsubsection{William Duer's Bankruptcy}

William Duer's leverage strategy---buying both securities and land with borrowed money---collapsed in March 1792:

\begin{itemize}
\item Unable to meet margin calls when securities market broke
\item Bankruptcy triggered the Panic of 1792
\item Imprisoned in New York's debtor's prison
\item Died in prison in 1799
\item Creditors eventually recovered approximately 50 cents on the dollar through long proceedings
\end{itemize}

\subsubsection{The General Pattern}

Most major land speculators of the 1780s failed financially:
\begin{itemize}
\item Those who purchased land expecting to pay in depreciated securities found their obligations had appreciated to par
\item Those who borrowed heavily to purchase land could not service debts when land sales failed to materialize
\item The depression of the mid-1790s destroyed land values
\item State laws against debtors were harsh, leading to imprisonment in many cases
\end{itemize}

\subsection{Effects on States}

\subsubsection{Winners from Assumption}

States with large unredeemed debts gained substantially:

\begin{enumerate}
\item \textbf{Massachusetts} ($\sim$\$4 million assumed):
\begin{itemize}
\item Had been heavily taxing citizens to service debt
\item Shays' Rebellion (1786--87) partly caused by these taxes
\item Assumption provided enormous fiscal relief
\item Citizens saw taxes decline as federal government took over debt service
\end{itemize}

\item \textbf{South Carolina} ($\sim$\$4 million assumed):
\begin{itemize}
\item Suffered heavy war damage
\item Planter class had large debts
\item Little capacity to service debt through taxation
\item Assumption was a major windfall
\end{itemize}

\item \textbf{Connecticut} ($\sim$\$1.6 million assumed):
\begin{itemize}
\item Had not serviced debt aggressively
\item Assumption relieved fiscal pressure
\end{itemize}
\end{enumerate}

\subsubsection{Losers from Assumption}

States that had already paid substantial debts felt aggrieved:

\begin{enumerate}
\item \textbf{Virginia} ($\sim$\$3.5 million assumed):
\begin{itemize}
\item Had already paid much more than \$3.5 million through state taxes
\item Citizens felt they were paying twice: once through past state taxes, again through future federal taxes to service other states' debts
\item Madison and Jefferson's opposition partly reflected Virginia's interests
\item Received compensation through capital location
\end{itemize}

\item \textbf{North Carolina} ($\sim$\$2.4 million assumed):
\begin{itemize}
\item Lost western lands (ceded to federal government 1790)
\item These lands became Tennessee (statehood 1796)
\item Ended lucrative land-for-debt schemes
\end{itemize}

\item \textbf{Georgia} ($\sim$\$2.2 million assumed):
\begin{itemize}
\item Federal government eventually took control of western lands
\item Ended Yazoo schemes
\item Led to decades of litigation
\end{itemize}
\end{enumerate}

\subsubsection{The Quid Pro Quo}

States accepting assumption generally had to:
\begin{itemize}
\item Cede western land claims to federal government
\item Accept federal supremacy in fiscal and debt matters
\item Give up land-for-debt schemes
\item Reduce their own taxing to avoid double taxation
\end{itemize}

This represented a fundamental shift in fiscal federalism from state to federal control.

\section{Long-Run Consequences}\label{sec:consequences}

\subsection{Establishment of Federal Fiscal Power}

Hamilton achieved his goal of consolidating fiscal power at the federal level. The funded debt created:

\begin{enumerate}
\item \textbf{Permanent revenue requirements:} Debt service required reliable federal revenues (tariffs, excises), establishing the principle of federal taxation

\item \textbf{Creditor constituency:} Bondholders had powerful financial interests in federal government strength and stability---exactly Hamilton's intention

\item \textbf{Precedent for federal supremacy:} Assumption established federal government's authority over major fiscal matters

\item \textbf{Financial capacity:} The funded debt and sound credit enabled future borrowing:
\begin{itemize}
\item Louisiana Purchase (1803): \$15 million borrowed from European banks
\item War of 1812: Substantial borrowing (though difficult given wartime conditions)
\item Later development projects and westward expansion
\end{itemize}
\end{enumerate}

\subsection{Development of Financial Markets}

Hamilton's funding system catalyzed American financial development:

\begin{enumerate}
\item \textbf{Securities markets:} Large-scale trading in government bonds emerged
\begin{itemize}
\item Professional securities dealers appeared in Philadelphia, New York, Boston
\item Trading volume exploded after funding
\item Led to formal stock exchanges
\end{itemize}

\item \textbf{The Buttonwood Agreement (1792):} Following the Duer panic, 24 brokers met under a buttonwood tree on Wall Street and established rules for securities trading---the origin of the New York Stock Exchange

\item \textbf{Bank of the United States (1791--1811):} Provided:
\begin{itemize}
\item Uniform national currency (bank notes)
\item Credit facilities for government and private sector
\item Professional financial management
\item Branch banking network
\end{itemize}

\item \textbf{Integration of capital markets:} Previously fragmented state-level markets became integrated nationally, improving capital allocation
\end{enumerate}

\subsection{Establishment of U.S. Creditworthiness}

The U.S. established an international reputation for honoring debts:

\begin{itemize}
\item By 1803, U.S. bonds traded at premium prices in London and Amsterdam
\item European investors (Dutch, British) eagerly purchased U.S. securities
\item Enabled Louisiana Purchase financing (borrowed from European banks)
\item Contrast with Latin American republics that defaulted repeatedly in 19th century
\item Tradition of debt service continued through 19th century (except brief Confederate default)
\end{itemize}

This creditworthiness became a valuable national asset, lowering borrowing costs for both government and private sector.

\subsection{Transformation of Western Land Policy}

Federal control of western lands meant:

\begin{enumerate}
\item \textbf{More orderly development:} Land surveyed systematically under Northwest Ordinance principles

\item \textbf{Federal revenues:} Land sales became a major revenue source for federal government (especially 1800--1860)

\item \textbf{End of state land-for-debt schemes:} No more arbitrage opportunities; lands sold through regular procedures

\item \textbf{Long-run development:} Eventually led to Homestead Act (1862) and systematic westward expansion

\item \textbf{Reduced corruption:} Federal administration (though far from perfect) was less corrupt than state land rushes of 1780s
\end{enumerate}

\subsection{Political Realignment: First Party System}

Hamilton's measures catalyzed the first American party system:

\textbf{Federalists:}
\begin{itemize}
\item Supported funded debt, national bank, strong federal government
\item Base: northern merchants, bondholders, commercial interests
\item Leaders: Hamilton, Washington, Adams, John Marshall
\item Dominant 1789--1800
\end{itemize}

\textbf{Democratic-Republicans:}
\begin{itemize}
\item Opposed assumption, suspicious of funded debt and national bank
\item Base: southern planters, yeoman farmers, agrarian interests
\item Leaders: Jefferson, Madison, Albert Gallatin
\item Dominant 1800--1824
\end{itemize}

This division reflected economic interests: the funded debt created a class whose wealth depended on federal government strength, exactly as Hamilton intended. The opposition represented those who saw federal power as threatening state sovereignty and agrarian interests.

\subsection{The Fate of the Speculator Class}

The great land speculators mostly failed, but their activities had lasting effects:

\begin{enumerate}
\item \textbf{Accelerated settlement:} Despite failures, speculators surveyed lands, built roads, established towns, and attracted settlers

\item \textbf{Risk-bearing:} Speculators bore enormous risks of frontier development; their failures transferred land to more successful hands

\item \textbf{Learned lessons:} The next generation of land developers (1800--1860) learned from the 1780s failures, using less leverage and more realistic pricing

\item \textbf{Creditor recoveries:} Many speculators' creditors eventually recovered partial amounts through lengthy legal proceedings

\item \textbf{Legal precedents:} The Yazoo cases established important contract law and constitutional principles (e.g., \textit{Fletcher v. Peck})
\end{enumerate}

\subsection{The Counterfactual: What if Hamilton Had Failed?}

Considering what might have occurred had assumption and funding failed:

\begin{enumerate}
\item \textbf{Debt depreciation:} Continental securities would likely have continued depreciating, possibly to zero

\item \textbf{Land speculation continuing:} States would have continued land-for-debt schemes until western lands were exhausted

\item \textbf{Weak federal government:} Without a funded debt and reliable revenues, federal government would have remained weak

\item \textbf{Possible fragmentation:} The United States might have fragmented into regional confederations (as occurred in Latin America after independence)

\item \textbf{No financial system:} Without funded debt and national bank, no basis for developed financial markets

\item \textbf{Foreign dependency:} Without domestic credit markets, greater dependence on European capital (as occurred in Latin America)

\item \textbf{Later development:} Economic development would have been slower without functioning capital markets
\end{enumerate}

The success of Hamilton's program thus appears crucial to American economic and political development.

\section{Theoretical Perspectives}\label{sec:theory}

\subsection{Ferguson's State-Building Interpretation}

E. James Ferguson interpreted assumption and funding as fundamentally about state-building rather than economic efficiency \citep{Ferguson1961}. Key insights:

\begin{enumerate}
\item \textbf{Creating a fiscal-military state:} Hamilton consciously modeled his system on British precedents---a funded debt, national bank, reliable revenues, and professional financial administration

\item \textbf{Binding creditor interests:} The funded debt created a powerful constituency whose wealth depended on federal strength, binding creditor interests to national government

\item \textbf{Centralizing power:} Assumption shifted fiscal power from states to federal government, establishing federal supremacy in taxation and borrowing

\item \textbf{Political coalition:} Creditors, merchants, and commercial interests became the base of the Federalist Party, supporting strong national government
\end{enumerate}

Ferguson's view emphasizes that Hamilton's goals were political as much as economic. Creating a class dependent on federal government was intentional state-building strategy.

\subsection{Rothbard's Special Interest Interpretation}

Murray Rothbard emphasized the conflict between competing special interests \citep{Rothbard1999, Rothbard2011}. His analysis stressed:

\begin{enumerate}
\item \textbf{Bondholder windfall:} Securities holders (often speculators who had purchased at deep discounts) reaped enormous gains at taxpayer expense

\item \textbf{Land speculator losses:} Many land speculators had legitimate businesses and took real risks; Hamilton's program destroyed their enterprises

\item \textbf{Taxpayer burden:} Ordinary citizens paid taxes to service the funded debt, transferring wealth to bondholders

\item \textbf{Special interest politics:} The Constitution itself and Hamilton's program served the interests of specific economic groups (merchants, creditors) rather than representing principled nationalism
\end{enumerate}

Rothbard's interpretation is more skeptical of Hamilton's motives, seeing special interest capture rather than visionary state-building.

\subsection{Modern Public Finance Perspective}

Contemporary economists might analyze these events through several lenses:

\subsubsection{Debt Restructuring and Incomplete Markets}

The situation resembled a sovereign debt restructuring:
\begin{itemize}
\item Old debt (Continental securities) traded at deep discounts due to uncertainty
\item New debt (funded federal bonds) offered lower yield but greater certainty
\item Restructuring resolved incomplete markets and commitment problems
\item Created liquid securities that served as money-like assets
\end{itemize}

\subsubsection{Fiscal Federalism}

The assumption debate involved optimal assignment of fiscal responsibilities:
\begin{itemize}
\item Which level of government should handle debt? Tax? Spend?
\item Trade-offs between local knowledge and economies of scale
\item Risk-sharing across states through federal assumption
\item Moral hazard: would assumption encourage future state profligacy?
\end{itemize}

\subsubsection{Time-Consistency Problems}

States faced commitment problems:
\begin{itemize}
\item Individual states could not credibly commit to repayment
\item Federal government with diversified revenue base could commit more credibly
\item Assumption solved a time-consistency problem in public debt
\end{itemize}

\subsubsection{Property Rights and Contract Theory}

Hamilton's program involved fundamental restructuring of property rights:
\begin{itemize}
\item Ended the implicit option to use securities as discounted land currency
\item Established securities as pure financial assets with reliable value
\item Created precedent for federal government honoring contracts
\item Established rule of law in financial markets
\end{itemize}

\subsection{Synthesis}

These perspectives are complementary rather than contradictory:
\begin{itemize}
\item Ferguson correctly identifies state-building and political coalition formation
\item Rothbard correctly identifies distributional conflicts and special interest politics
\item Modern public finance correctly identifies commitment problems and optimal fiscal federalism
\end{itemize}

Hamilton's program simultaneously:
\begin{enumerate}
\item Built state capacity and federal power (Ferguson)
\item Redistributed wealth from land speculators and taxpayers to bondholders (Rothbard)
\item Solved commitment problems and created functioning financial markets (public finance)
\end{enumerate}

All three occurred; the question is which effects were intended, which were incidental, and which normative framework one adopts for evaluation.

\section{Conclusion}\label{sec:conclusion}

The conflict between land speculators and public creditors in the 1780s--90s reveals much about American state-building and fiscal development. Several conclusions emerge:

\begin{enumerate}
\item \textbf{Fundamental incompatibility:} The two uses of depreciated securities---as land currency versus financial assets---were fundamentally incompatible. Hamilton's funding resolved this conflict decisively in favor of creditors.

\item \textbf{Distributional consequences:} Public creditors gained enormously; land speculators lost enormously. These weren't small effects but rather complete financial ruin for major figures like Morris, Duer, and Gorham, versus vast windfall gains for bondholders.

\item \textbf{State-building success:} Whatever one's normative evaluation, Hamilton succeeded in his goals. He created:
\begin{itemize}
\item A funded national debt
\item Federal fiscal capacity
\item A creditor constituency supporting federal power
\item Sound public credit
\item The foundation for financial markets
\end{itemize}

\item \textbf{Path dependence:} The choices made in 1790--91 set the United States on a path of strong federal fiscal power, reliable debt service, and developed financial markets. The contrast with Latin America (which followed different paths) suggests these choices mattered greatly.

\item \textbf{Theoretical insights:} The episode illustrates:
\begin{itemize}
\item How fiscal policy can be used for state-building (Ferguson)
\item The importance of distributional conflicts in policy (Rothbard)
\item Solutions to commitment problems and fiscal federalism (public finance)
\end{itemize}

\item \textbf{Lessons for modern fiscal policy:} Several lessons carry forward:
\begin{itemize}
\item Credibility matters: Hamilton's commitment to debt service established valuable reputation
\item Distributional effects matter: Funding enriched creditors but bankrupted speculators
\item Fiscal institutions matter: Choice of federal vs. state fiscal responsibility has long-run consequences
\item Financial development requires reliable government debt
\item Political economy constraints bind: Even Hamilton needed complex vote-trading to pass assumption
\end{itemize}
\end{enumerate}

The land speculation versus funding conflict of 1783--95 thus represents a formative episode in American fiscal history, with consequences extending far beyond the immediate distributional effects to shape the long-run development of American state capacity, financial markets, and political economy.

\bibliographystyle{plainnat}
\bibliography{fiscal_history}

\end{document}
